\documentclass[exercice,a4]{thedocument}%
\usepackage{pstricks}%
\usepackage{tabularx}%
\usepackage{textcomp}%
\startdatakeys
    \source{25GENMATPO3}
    \BO{2026}
    \cycle{4}
    \competencesDuSocle{CA,RA}
    \objectifsApprentissage{C4A1aa,C4A2ad,C4B2cb,C4D2aa,C4D3ac}
    \motsCles{trigonométrie,vitesse,échelle,coordonnées géographiques}
\enddatakeys
\begin{document}%
    
    \def\degres{\textsuperscript{\textdegree}}%

\medskip

Cet exercice est un questionnaire à choix multiple (QCM).

Pour chaque question, quatre réponses sont proposées. \textbf{Une seule réponse est exacte}.

Recopier sur la copie le numéro de la question et la réponse choisie. Aucune justification n'est demandée.

\medskip

\textbf{Question 1}

La formule qui permet d’obtenir la pointure de chaussure $p$ à partir de la longueur $L$ du pied, exprimée en centimètres, est  $p = 1,5(L + 1)$.

Quelle est la pointure d’un pied de longueur 25 cm ?

\begin{center}
\begin{tabularx}{\linewidth}{|*{4}{>{\centering \arraybackslash}X|}}\hline
37,5& 38& 38,5&39\\ \hline
\end{tabularx}
\end{center}

\textbf{Question 2}

\begin{minipage}{0.6\linewidth}
Le triangle MUR, rectangle en U, qui est représenté ci-contre n’est pas en vraie grandeur.

Quelle est la mesure de l'angle $\widehat{\text{RMU}}$ arrondie au degré ?
\end{minipage}\hfill
\begin{minipage}{0.35\linewidth}
\psset{unit=1cm}
\begin{pspicture}(5,3.3)
\pspolygon(0.3,0.3)(3.8,0.3)(3.8,2.8)
\psframe(3.8,0.3)(3.5,0.6)
\psarc(0.3,0.3){0.9}{0}{37}
\uput[dl](0.3,0.3){M} \uput[dr](3.8,0.3){U}\uput[ur](3.8,2.8){R}\uput[ul](1.8,1.6){9,7 cm}\uput[r](3.8,1.55){7,2 cm}\uput[ur](1.4,0.5){?}
\end{pspicture}
\end{minipage}

\begin{center}
\begin{tabularx}{\linewidth}{|*{4}{>{\centering \arraybackslash}X|}}\hline
$37\degres$& $42\degres$&$48\degres$&$53\degres$\\ \hline
\end{tabularx}
\end{center}

\textbf{Question 3}

Un coureur a remporté la course du 100~m en 10 secondes exactement.

À quelle vitesse moyenne en km/h a-t-il parcouru cette distance ?

\begin{center}
\begin{tabularx}{\linewidth}{|*{4}{>{\centering \arraybackslash}X|}}\hline
10 km/h &36 km/h& 44 km/h&60 km/h\\ \hline
\end{tabularx}
\end{center}

\textbf{Question 4}

Roxane a acheté le modèle réduit d’un tableau célèbre à l’échelle $\dfrac{1}{20}$.

La largeur du modèle réduit est 7 cm. Quelle est la largeur réelle du tableau ?

\begin{center}
\begin{tabularx}{\linewidth}{|*{4}{>{\centering \arraybackslash}X|}}\hline
 14 cm &27 cm& 35 cm&140 cm\\ \hline
\end{tabularx}
\end{center}

\textbf{Question 5}

\begin{minipage}{0.45\linewidth}
Trois points nommés A, B, et C sont positionnés sur la sphère ci-contre.
\end{minipage}\hfill
\begin{minipage}{0.53\linewidth}
\psset{unit=1cm}
\begin{pspicture}(8,4.7)
\pscircle(2.5,2.4){2}
\psdots[dotstyle=+,dotangle=45,dotscale=2.2](1.48,3.4)(1.85,0.75)(3.04,3.35)
\uput[ur](1.48,3.4){\small A}\uput[ur](3.04,3.35){\small C}\uput[ur](1.85,0.75){\small B}
\psellipticarc(2.5,2.4)(2,0.4){-180}{0}
\psellipticarc[linestyle=dashed](2.5,2.4)(2,0.4){0}{180}
\psellipticarc(2.5,1.6)(1.85,0.4){-180}{0}
\psellipticarc[linestyle=dashed](2.5,1.6)(1.85,0.4){0}{180}
\psellipticarc(2.5,1.)(1.38,0.3){-180}{0}
\psellipticarc[linestyle=dashed](2.5,1)(1.38,0.3){0}{180}
\psellipticarc(2.5,3.6)(1.58,0.3){-180}{0}
\psellipticarc[linestyle=dashed](2.5,3.6)(1.58,0.3){0}{180}
\psellipticarc(2.5,2.4)(1.2,2){90}{270}
\psellipticarc(2.5,2.4)(0.6,2){-90}{90}
\psellipticarc(2.5,2.4)(1.2,2){-90}{90}
\psellipticarc(2.5,2.4)(0.3,2){90}{-90}
\rput(4.4,3.6){\footnotesize $65\degres$}\rput(4.7,2.4){\footnotesize $0\degres$}\rput(4.6,1.5){\footnotesize $40\degres$}\rput(4.2,0.8){\footnotesize $65\degres$}
\rput(2.5,4.6){\footnotesize N}\rput(2.5,0.2){\footnotesize S}\rput(0.2,2.4){\footnotesize O}
\rput(5,2.4){\footnotesize E}
\uput[r](5.4,3.6){N : Nord} \uput[r](5.4,2.8){S : Sud}
\uput[r](5.4,2){O : Ouest} \uput[r](5.4,1.2){E : Est}
\rput{90}(1.1,2.4){\footnotesize $30\degres$}\rput{90}(2,2.2){\footnotesize $0\degres$}
\rput{90}(2.8,2.3){\footnotesize $30\degres$}\rput{90}(3.5,2.4){\footnotesize $50\degres$}
\end{pspicture}

\end{minipage}

Quel point a pour coordonnées ($65\degres$ N~;~$30\degres$ O) ?

\begin{center}
\begin{tabularx}{\linewidth}{|*{4}{>{\centering \arraybackslash}X|}}\hline
A &B& C& Aucun\\ \hline
\end{tabularx}
\end{center}

\bigskip
\showthesource
\end{document}
